% Mathematics Paper -- an amsart-inspired research paper on A4.
% The page design (grid, Libertinus type, theorem styles, running heads,
% accent colour) lives in mathpaper.sty; this file holds the content.
% Edit the title, authors, MSC codes and addresses in the metadata block.
% Build with: latexmk -pdf main.tex
\documentclass[11pt, twoside]{article}
\usepackage{mathpaper}

\usepackage{tikz}
\usetikzlibrary{arrows.meta, decorations.pathreplacing}
\usepackage{tikz-cd}
\tikzcdset{arrow style=tikz, diagrams={>={Straight Barb[scale=0.8]}}}

% Numeric references in brackets, small, with initials as in AMS journals.
\usepackage[backend=biber, style=numeric-comp, sorting=nyt, giveninits=true,
  maxbibnames=5, doi=false, url=false, isbn=false, eprint=false]{biblatex}
\addbibresource{refs.bib}
\renewcommand*{\bibfont}{\small}
\DeclareFieldFormat[article]{title}{\mkbibemph{#1}}
\DeclareFieldFormat[book]{title}{\mkbibemph{#1}}
\DeclareFieldFormat{journaltitle}{#1}
\renewbibmacro{in:}{}

% Links in the accent colour, without boxes.
\usepackage[colorlinks, allcolors=accent]{hyperref}
\usepackage[capitalise, nameinlink]{cleveref}
% Theorem A is unnumbered, so it is referenced by name.
\newcommand{\thmA}{\hyperref[thm:A]{Theorem~A}}

% Notation used throughout; change it in one place.
\newcommand{\R}{\mathbb{R}}
\def\Bbbk{\mathbf{k}}
\newcommand{\vect}{\mathbf{vect}_{\Bbbk}}
\newcommand{\Pers}{\mathbf{Pers}}
\newcommand{\sh}[2]{#1[#2]}
\DeclareMathOperator{\im}{im}
\DeclareMathOperator{\coker}{coker}
\DeclareMathOperator{\Barc}{Bar}
\newcommand{\dI}{d_{\mathrm{I}}}
\newcommand{\dB}{d_{\mathrm{B}}}

% ---- Metadata --------------------------------------------------------------
\title{Mayer--Vietoris Sequences\\and Interleavings of Persistence Modules}
\shorttitle{Mayer--Vietoris sequences and interleavings}
\author{Agnieszka Wróbel \and Tomás Echeverría}
\renewcommand{\and}{\enspace and\enspace}
\shortauthors{A. Wróbel and T. Echeverría}
\date{Received 3 March 2026; revised 14 August 2026}
\subjclass{Primary 55N31; Secondary 18G80, 62R40.}
\keywords{persistence modules, interleaving distance, Mayer--Vietoris
  sequence, barcodes, covers of filtered spaces.}
\address{Agnieszka Wróbel}{Institute of Mathematics, Jagiellonian
  University, ul.\ Łojasiewicza 6, 30-348 Kraków, Poland}{a.wrobel@example.org}
\address{Tomás Echeverría}{Departamento de Matemáticas, Universidad de
  Chile, Las Palmeras 3425, Ñuñoa, Santiago, Chile}{t.echeverria@example.org}

\begin{abstract}
Given a filtered space covered by two filtered subspaces, the persistent
homology of the whole is computed from the pieces by a long exact
Mayer--Vietoris sequence, but in general the pieces do not determine it up to
isomorphism. We show that they determine it up to an $\varepsilon$-interleaving
whenever the image of the connecting morphism is $\varepsilon$-trivial, and
that this hypothesis holds as soon as every bar of the intersection in the
adjacent degree is shorter than $\varepsilon$. By the isometry theorem the
same bound holds for the bottleneck distance between barcodes. We give a
quantitative version for covers that are only approximately excisive and
illustrate the bound on a sampled annulus.
\end{abstract}

\begin{document}
\maketitle

\section{Introduction}
\label{sec:intro}

Persistent homology assigns to a filtered space $X = (X_t)_{t\in\R}$ the
family of vector spaces $H_n(X_t)$ together with the maps induced by the
inclusions $X_s \subseteq X_t$. Such a family is a \emph{persistence module},
and when each space is finite-dimensional it decomposes into interval modules,
recorded as a barcode~\cite{crawley-boevey2015,zomorodian2005}. Barcodes are
stable: nearby filtrations have nearby barcodes, and the interleaving distance
between modules equals the bottleneck distance between their
barcodes~\cite{chazal2009,lesnick2015,bauer2015}.

Computations rarely see a whole space at once. In distributed and multiscale
settings one computes the persistent homology of pieces $A$ and $B$ with
$X = A \cup B$ and must recover that of $X$. Levelwise, the Mayer--Vietoris
sequence
\begin{equation}
\label{eq:mv-level}
\cdots \to H_n(A_t\cap B_t) \to H_n(A_t)\oplus H_n(B_t) \to H_n(X_t)
  \xrightarrow{\ \delta\ } H_{n-1}(A_t\cap B_t) \to \cdots
\end{equation}
is natural in $t$, hence a long exact sequence of persistence
modules~\cite{dicuroya2023}. Exactness does not give $H_n(X)$, however: an
extension problem remains, and in persistence it is unsolvable in general.
Our observation is that an extension by a module whose structure maps die
quickly is close to split, and that ``die quickly'' can be read off the
barcode of the intersection.

\begin{maintheorem}
\label{thm:A}
Let $X = A\cup B$ be a filtered space covered by the interiors of the
filtered subspaces $A$ and $B$, and let $Q_n$ be the cokernel of
$H_n(A\cap B) \to H_n(A)\oplus H_n(B)$. If every bar of $H_{n-1}(A\cap B)$ has
length at most $\varepsilon$, then $\dI\bigl(H_n(X), Q_n\bigr) \le \varepsilon$
and $\dB\bigl(\Barc H_n(X), \Barc Q_n\bigr) \le \varepsilon$.
\end{maintheorem}

The theorem is sharp: for every $\varepsilon > 0$ there are covers of a
filtered circle for which both distances equal $\varepsilon$
(\cref{ex:annulus}). \Cref{sec:prelim} recalls interleavings,
\cref{sec:extensions} proves the algebraic lemma behind the theorem, and
\cref{sec:covers} applies it to covers and to covers that are only
approximately excisive.

\section{Preliminaries}
\label{sec:prelim}

Fix a field $\Bbbk$ and write $\vect$ for the category of finite-dimensional
vector spaces over $\Bbbk$. We regard $\R$ as a poset category.

\begin{definition}
A \emph{persistence module} is a functor $M\colon \R \to \vect$. We write $M_t$
for its value at $t$ and $\varphi_M^{s,t}\colon M_s \to M_t$ for its
\emph{structure map} when $s \le t$. Morphisms are natural transformations,
and $\Pers$ denotes the resulting abelian category.
\end{definition}

For $\varepsilon \ge 0$ the \emph{shift} $\sh{M}{\varepsilon}$ is the module
with $\sh{M}{\varepsilon}_t = M_{t+\varepsilon}$, and the structure maps of
length $\varepsilon$ assemble into a morphism
$\sigma_M^\varepsilon\colon M \to \sh{M}{\varepsilon}$. A morphism
$f\colon M\to N$ shifts to $\sh{f}{\varepsilon}$ in the evident way.

\begin{definition}
\label{def:interleaving}
Modules $M$ and $N$ are \emph{$\varepsilon$-interleaved} if there are
morphisms $f\colon M \to \sh{N}{\varepsilon}$ and
$g\colon N \to \sh{M}{\varepsilon}$ such that the triangles
\begin{equation}
\label{eq:interleaving}
\begin{tikzcd}[column sep=small, row sep=large]
M \arrow[rr, "\sigma_M^{2\varepsilon}"] \arrow[dr, "f"']
  & & \sh{M}{2\varepsilon} \\
& \sh{N}{\varepsilon} \arrow[ur, "\sh{g}{\varepsilon}"'] &
\end{tikzcd}
\qquad\qquad
\begin{tikzcd}[column sep=small, row sep=large]
& \sh{M}{\varepsilon} \arrow[dr, "\sh{f}{\varepsilon}"] & \\
N \arrow[rr, "\sigma_N^{2\varepsilon}"'] \arrow[ur, "g"]
  & & \sh{N}{2\varepsilon}
\end{tikzcd}
\end{equation}
commute. The \emph{interleaving distance} is
\[
\dI(M,N) = \inf\{\varepsilon \ge 0 : M \text{ and } N \text{ are
$\varepsilon$-interleaved}\},
\]
with the convention $\inf\emptyset = \infty$.
\end{definition}

The interleaving distance is an extended pseudometric on $\Pers$; the
triangle inequality follows by composing interleavings~\cite{bubenik2014}.
Every pointwise finite-dimensional module is a direct sum of interval modules
$\Bbbk_I$, and the multiset of intervals is its \emph{barcode}
$\Barc M$~\cite{crawley-boevey2015}.

\begin{theorem}[Isometry theorem~\cite{lesnick2015,bauer2015}]
\label{thm:isometry}
For pointwise finite-dimensional modules $M$ and $N$,
\begin{equation}
\label{eq:isometry}
\dI(M,N) = \dB(\Barc M, \Barc N).
\end{equation}
\end{theorem}

\begin{definition}
A module $K$ is \emph{$\varepsilon$-trivial} if $\sigma_K^\varepsilon = 0$,
that is, if every structure map of length $\varepsilon$ vanishes.
\end{definition}

\begin{remark}
\label{rem:trivial}
A pointwise finite-dimensional module is $\varepsilon$-trivial if and only if
each of its bars has length at most $\varepsilon$. Submodules and quotients of
$\varepsilon$-trivial modules are $\varepsilon$-trivial, because the
structure maps of a submodule are restrictions and those of a quotient are
induced maps.
\end{remark}

\section{Extensions by trivial modules}
\label{sec:extensions}

The following lemma carries the whole argument. It says that an extension by
an $\varepsilon$-trivial module is $\varepsilon$-interleaved with its
submodule, whatever the extension class.

\begin{lemma}
\label{lem:extension}
Let $0 \to Q \xrightarrow{\iota} M \xrightarrow{\pi} K \to 0$ be a short exact
sequence in $\Pers$ with $K$ $\varepsilon$-trivial. Then there is a unique
morphism $\psi\colon M \to \sh{Q}{\varepsilon}$ with
$\sh{\iota}{\varepsilon}\circ\psi = \sigma_M^\varepsilon$, and $\iota$ and
$\psi$ form an $\varepsilon$-interleaving. In particular
$\dI(Q, M) \le \varepsilon$.
\end{lemma}

\begin{proof}
Naturality of $\pi$ gives $\sh{\pi}{\varepsilon}\circ\sigma_M^\varepsilon =
\sigma_K^\varepsilon\circ\pi = 0$, so $\sigma_M^\varepsilon$ lands in the
kernel of $\sh{\pi}{\varepsilon}$, which is the image of the monomorphism
$\sh{\iota}{\varepsilon}$. This yields $\psi$ and its uniqueness, and the
diagram
\begin{equation}
\label{eq:extension}
\begin{tikzcd}
0 \arrow[r] & Q \arrow[r, "\iota"] \arrow[d, "\sigma_Q^\varepsilon"']
  & M \arrow[r, "\pi"] \arrow[d, "\sigma_M^\varepsilon"]
    \arrow[dl, dashed, "\psi" description]
  & K \arrow[r] \arrow[d, "0"] & 0 \\
0 \arrow[r] & \sh{Q}{\varepsilon} \arrow[r, "\sh{\iota}{\varepsilon}"']
  & \sh{M}{\varepsilon} \arrow[r, "\sh{\pi}{\varepsilon}"']
  & \sh{K}{\varepsilon} \arrow[r] & 0
\end{tikzcd}
\end{equation}
commutes. Put $f = \sh{\iota}{\varepsilon}\circ\sigma_Q^\varepsilon\colon
Q \to \sh{M}{\varepsilon}$ and $g = \psi$. The two triangles of
\eqref{eq:interleaving} read
\[
\sh{\psi}{\varepsilon}\circ f = \sigma_Q^{2\varepsilon}
\quad\text{and}\quad
\sh{f}{\varepsilon}\circ\psi = \sigma_M^{2\varepsilon}.
\]
The second follows from $\sh{\iota}{\varepsilon}\circ\psi =
\sigma_M^\varepsilon$ and naturality of $\sigma$. Composing the first with
$\sh{\iota}{2\varepsilon}$ reduces it to the same identity, and
$\sh{\iota}{2\varepsilon}$ is injective, so it may be cancelled.
\end{proof}

\begin{remark}
The constant in \cref{lem:extension} is sharp. For $M = \Bbbk_{[0,\infty)}$
and its submodule $Q = \Bbbk_{[\varepsilon,\infty)}$ the quotient is
$K = \Bbbk_{[0,\varepsilon)}$, which is $\varepsilon$-trivial, and
$\dI(Q,M) = \varepsilon$ because the infinite bars must be matched to each
other.
\end{remark}

\section{Covers of filtered spaces}
\label{sec:covers}

A \emph{filtered space} is a functor $X\colon \R \to \mathbf{Top}$ whose
structure maps are inclusions, and a \emph{cover} of $X$ is a pair of filtered
subspaces with $A_t \cup B_t = X_t$ for all $t$. The cover is
\emph{excisive} if $\{\mathrm{int}\,A_t, \mathrm{int}\,B_t\}$ covers $X_t$ for
every $t$; then \eqref{eq:mv-level} is exact at each $t$ and, being natural,
is a long exact sequence in $\Pers$.

\begin{proof}[Proof of \thmA]
Cutting the long exact sequence at $H_n(X)$ gives a short exact sequence
\begin{equation}
\label{eq:cut}
0 \to Q_n \to H_n(X) \xrightarrow{\ \delta\ } \im\delta \to 0,
\end{equation}
where $\im\delta \subseteq H_{n-1}(A\cap B)$. By hypothesis and
\cref{rem:trivial}, $H_{n-1}(A\cap B)$ and hence $\im\delta$ are
$\varepsilon$-trivial. \Cref{lem:extension} gives
$\dI(H_n(X), Q_n) \le \varepsilon$, and \cref{thm:isometry} turns this into
the bound on barcodes.
\end{proof}

\begin{figure}[t]
\centering
\begin{tikzpicture}[scale=0.92, line cap=round]
  % Left: an annulus covered by two overlapping arcs A (left) and B (right).
  \begin{scope}
    \fill[black!6, even odd rule] (0,0) circle (1.55) (0,0) circle (0.85);
    \draw[black!55, line width=.5pt] (0,0) circle (1.55) (0,0) circle (0.85);
    \draw[accent, line width=2.4pt, opacity=.85]
      (100:1.2) arc[start angle=100, end angle=260, radius=1.2];
    \draw[black!70, line width=2.4pt, opacity=.55]
      (-80:1.2) arc[start angle=-80, end angle=80, radius=1.2];
    \node[accent] at (180:1.95) {$A$};
    \node[black!75] at (0:1.95) {$B$};
    \node[font=\footnotesize, black!60] at (0,-2.05) {cover of $X_t$};
  \end{scope}
  % Right: barcodes with the bars of the intersection kept short.
  \begin{scope}[xshift=4.3cm, yshift=-1.25cm]
    \draw[black!45, -{Straight Barb[scale=.7]}] (0,0) -- (5.2,0)
      node[right, font=\footnotesize, black!60] {$t$};
    \foreach \x/\l in {0/0, 1/1, 2/2, 3/3, 4/4}
      \draw[black!45] (\x,0) -- (\x,-.08) node[below, font=\scriptsize] {\l};
    \node[left, font=\footnotesize] at (-.15,2.45) {$H_0(A\cap B)$};
    \draw[black!60, line width=1.6pt] (.2,2.45) -- (4.9,2.45);
    \draw[black!60, line width=1.6pt] (.4,2.15) -- (1.05,2.15);
    \node[left, font=\footnotesize] at (-.15,1.35) {$H_1(X)$};
    \draw[accent, line width=1.6pt] (1.05,1.35) -- (4.4,1.35);
    \node[left, font=\footnotesize] at (-.15,.6) {$Q_1$};
    \draw[black!35, line width=1.6pt] (1.7,.6) -- (4.4,.6);
    \draw[decorate, decoration={brace, amplitude=3pt, mirror}, black!60]
      (.4,1.98) -- (1.05,1.98) node[midway, below=2pt, font=\scriptsize]
      {$\le\varepsilon$};
  \end{scope}
\end{tikzpicture}
\caption{A sampled annulus covered by two arcs (left) and the barcodes in
\cref{ex:annulus} (right). The short bar of $H_0(A\cap B)$ controls the
difference between the bar of $H_1(X)$ and its approximation from the pieces.}
\label{fig:annulus}
\end{figure}

\begin{example}
\label{ex:annulus}
Let $X$ be the Vietoris--Rips filtration of a dense sample of an annulus,
covered by the filtrations $A$ and $B$ of two overlapping arcs as in
\cref{fig:annulus}. The intersection $A_t\cap B_t$ has two components for
small $t$, one of which merges after time $\varepsilon$; the class in
$H_1(X)$ that goes around the annulus is born when the second overlap
appears and maps under $\delta$ onto the short bar. The cokernel $Q_1$ sees
the loop only once the short bar has died, so the left endpoints of the
two bars differ by exactly $\varepsilon$, and the bound of \thmA{}
is attained.
\end{example}

Covers built from data, as in the distributed algorithms
of~\cite{yoon2023,casas2021}, are seldom excisive at every scale: two patches of a
point cloud overlap only after their Rips complexes have grown. The next
result measures this defect and pays for it in the bound.

\begin{definition}
A cover is \emph{$\eta$-excisive} if for all $t$ the inclusion of the
singular chains of $X_t$ subordinate to $\{A_t, B_t\}$ into all singular
chains of $X_t$ becomes a homology isomorphism after composing with the
inclusion into $X_{t+\eta}$, in every degree.
\end{definition}

\begin{proposition}
\label{prop:approx}
Suppose that the cover $X = A\cup B$ is $\eta$-excisive and that every bar of
$H_{n-1}(A\cap B)$ has length at most $\varepsilon$. Then
\begin{equation}
\label{eq:approx}
\dI\bigl(H_n(X), Q_n\bigr) \le \varepsilon + \eta .
\end{equation}
\end{proposition}

\begin{proof}
Let $C$ be the filtered complex of chains subordinate to the cover. Its
homology fits into an exact Mayer--Vietoris sequence for every $t$, so the
proof of \thmA{} gives $\dI(H_n(C), Q_n) \le \varepsilon$. The
hypothesis says that the inclusion $H_n(C) \to H_n(X)$ and the map
$H_n(X) \to \sh{H_n(C)}{\eta}$ defined through $X_{t+\eta}$ form an
$\eta$-interleaving, so $\dI(H_n(C), H_n(X)) \le \eta$. The triangle
inequality for $\dI$ completes the proof.
\end{proof}

\begin{corollary}
\label{cor:local}
Under the hypotheses of \cref{prop:approx},
$\dB\bigl(\Barc H_n(X), \Barc Q_n\bigr) \le \varepsilon + \eta$. In particular
the barcode of $X$ in degree $n$ is determined up to $\varepsilon + \eta$ by
$A$, $B$ and $A\cap B$ alone.
\end{corollary}

\subsection*{Acknowledgements}
We thank the referee for a careful reading and for suggesting the sharper
form of \cref{lem:extension}. The first author was supported by a fictional
grant of the National Science Centre.

\printbibliography
\printaddresses

\end{document}
